%================================================================================
%       Safety Critical Systems Club - Data Safety Initiative Working Group
%================================================================================
%                       DDDD    SSSS  IIIII  W   W   GGGG
%                       D   D  S        I    W   W  G   
%                       D   D   SSS     I    W W W  G  GG
%                       D   D      S    I    WW WW  G   G
%                       DDDD   SSSS   IIIII  W   W   GGG
%================================================================================
%               Data Safety Guidance Document - LaTeX Source File
%================================================================================
%
% Description:
%   Activities and Tailoring section.
%
%================================================================================
\chapter{Tailoring (Informative)}\index{Tailoring|textbf}\label{bkm:activitiestailoring}
\todo{Describe relationship to activities in Part 1, and cross-reference both ways.}
\dsiwgSectionQuote
  {It is a capital mistake to theorise before one has data.}
  {Sherlock Holmes - ``A Study in Scarlet'' (Sir Arthur Conan Doyle)}


\section{Phase 1: Establish Context}
The level of
\index{Stakeholder}\gls{stakeholder}
interface control\index{Interface!organizational} required will depend heavily on the number of \index{Stakeholder}\glspl{stakeholder}, complexity of their interactions, and the contractual controls already in place. Many programmes already require \glspl{icd}\index{Interface!Control Document} to be developed for systems or equipment. The level of detail required in other programme interface control plans may be used as a guide for the requirements of data safety interface control. 

The guidance in this document is general in nature, and so it is likely to need \gls{tailoring}\index{Tailoring} to align with the organization's attitude to risk, their existing processes and the relevant sector's regulatory environment. The provided \gls{odr}, or a version customised to suit the organization's needs, may therefore provide a structured approach to \cbstart data safety management\cbend.
The \gls{odr} assessment can be conducted at the product line or the individual product level, as appropriate for the organization. It is generally not recommended to conduct the \gls{odr} without the context of a system type.

The \gls{odr} assessment is likely to be most useful for organizations that do not have significant safety engineering experience and that are operating in industrial sectors that are not strongly regulated.
organizations with considerable experience in the development of safety-critical systems in heavily regulated environments
may also find it useful for defining data context and for augmenting any existing safety standards which do not explicitly consider data.
If this assessment is not conducted, some high-level qualitative estimate of risk may still be required (e.g., to support process \gls{tailoring}\index{Tailoring}); likewise, there will also be a need to identify key \index{Stakeholder}\glspl{stakeholder} and necessary approvers.

The Data Safety Culture questionnaire is likely to be of most use for low-risk (i.e., ODR1) systems. Developers of higher risk systems will typically have existing processes to develop, maintain and monitor safety cultures, although the data-oriented questionnaire could help inform those existing processes.

Including data safety within a project safety management plan is recommended for complex or highly safety-critical systems. In this case the structure of the safety management plan may be maintained, with the \index{Safety Assessment!Data}\gls{safety assessment} process being tailored to align to the overall \cbstart data safety management\cbend\ process.

The data maturity questionnaire is likely to be of most use when new organizational relationships are being formed. It may offer little value in situations where both organizations are familiar with each other, they have worked on data-related projects together before and there are suitable audit / review arrangements in place.

\index{Artefact, Data}\Glspl{data artefact} may be defined at a number of levels. A \gls{data artefact} associated with a medical system could be described as ``patient data''. Alternatively, this could be split into smaller parts (e.g., ``blood group''). Generally, the highest possible level consistent with the system description should be used; this prevents an excessively long list of \glspl{data artefact} being developed. If necessary, where further detail is needed those \glspl{data artefact} can be refined as part of an iterative process focused on key issues.

Not every data category\index{Data!Category|} will be relevant to every system. Furthermore, for low-risk (ODR1) systems it may be sufficient to simply consider the groupings of categories (e.g., ``context'', ``implementation'', etc.). Conversely, high-risk (ODR4) systems might need to consider every category\index{Category!Data}, even if this results in a conclusion that a specific category\index{Category!Data} is not relevant for the system in question.

A function-based approach to identifying \index{Artefact, Data}\glspl{data artefact} is likely to be enabled by design processes that also adopt a function-based perspective. If \gls{information} from a function-based perspective is readily available, then it should be used to support the identification of \index{Artefact, Data}\glspl{data artefact}. If this \gls{information} is not readily available, it is recommended that it be generated for medium and high-risk systems (i.e., ODR2 to ODR4, inclusive).

Considering data across the system lifecycle\index{Lifecycle!System} is a relatively simple activity, which is applicable to all systems (i.e., ODR1 to ODR4, inclusive).

\section{Phase 2: Identify Risks}
The general, historical perspective review is a simple activity that does not require significant resources. It is recommended for all systems, regardless of risk level.

When conducting a bottom-up approach, it may not be appropriate to explicitly consider every possible property for every single artefact. For low-risk (ODR1) and medium-risk (ODR2 / ODR3) systems some form of \gls{tailoring}\index{Tailoring} may be expected. This might, for example, take the form of pre-selecting the properties that are most relevant, or limiting the layer of abstraction at which the system is considered.

\Gls{tailoring}\index{Tailoring} of the bottom-up approach may also be appropriate for some high-risk (ODR4) systems, but in this case an explicit argument that the \gls{tailoring}\index{Tailoring} has not adversely affected system safety would be expected. Furthermore, the risk identification process for high-risk (ODR4) systems is expected to be highly structured. To support this, a number of data-related \gls{hazop}\index{HAZOP} guidewords are presented in \autoref{bkm:guidewords}.

The top-down and bottom-up approaches provide different perspectives on data-related risks. For low-risk (ODR1) systems it may be appropriate to consider just one of these perspectives. Both perspectives would be expected to be considered (to some degree) for medium-risk (ODR2 / ODR3) and high-risk (ODR4) systems.

\section{Phase 3: Analyse Risks}\label{bkm:activities:analyse:tailoring}
\glspl{dsal}\index{Assurance Level!Data} can be applied at different levels and to different constructs. For example, in the case of simple, low-risk systems it may be appropriate to apply a single \gls{dsal} to an entire system. Alternatively, it may be appropriate to apply \glspl{dsal} to sub-systems, for example, to match the level at which other \gls{integrity}, or assurance, levels have been determined. 

Another option is to apply \glspl{dsal}\index{Assurance Level!Data} to \glspl{data artefact} rather than directly to risks. This approach has the advantage that \index{Treatment!Risk}\glspl{treatment} are often related to \glspl{data artefact}; it can work well where there is a simple relationship between \glspl{data artefact} and risks.

\section{Phase 4: Evaluate and \Glsfmttext{treat} Risks}
Records of the discussions that occur as part of risk evaluation should be kept. For low-risk (ODR1) systems, this may be in the form of a brief memo. For high-risk (ODR4) systems, a detailed, structured record, which is placed under formal control, may be required; in this case these discussions may be recorded as part of the system's \gls{hazard log} or as part of a data safety management plan.

As outlined in
\hyperref[bkm:activities:analyse:tailoring]{Section}~\ref{bkm:activities:analyse:tailoring}, %Clunky, but \autoref inserted "subsubsection"
\glspl{dsal}\index{Assurance Level!Data} can be applied at varying levels of abstraction. For small-scale or low-risk (ODR1) systems it may be appropriate to consider \index{Treatment!Risk}\glspl{treatment} at higher levels of system abstraction. For example, this could involve applying a single \gls{dsal} to a sub-system or even to the system in its entirety and implementing risk \index{Treatment!Risk}\gls{treatment} techniques at that level. The latter approach could be appropriate where the data in the system interacts in complex ways and the associated safety risk does not warrant a detailed investigation of these interactions.

A significant amount of \gls{tailoring}\index{Tailoring} is implicit in the way the tables of methods and approaches are constructed. \hyperref[bkm:activities:analyse:tailoring]{Section}~\ref{bkm:activities:analyse:tailoring} recommends methods and approaches as ways of maintaining a required data property at a given \gls{dsal}, but do not mandate those methods.
